% Module 9: Quantum Simulations
% Estimated slides: 25 | Duration: ~2 hours
% Key topics: Feynman's vision, quantum chemistry simulation, VQE for molecules,
%             material science, drug discovery, financial simulation
% Labs: H2 molecule VQE in Qiskit Nature; LiH ground state

%%%%%%%%%%%%%%%%%%%%%%%%%%%%%%%%%%%%%%%%%%%%%%%%%%%%%%%%%%
\begin{frame}[fragile]\frametitle{}
\begin{center}
{\Large Module 9: Quantum Simulations}
\end{center}
\end{frame}

%%%%%%%%%%%%%%%%%%%%%%%%%%%%%%%%%%%%%%%%%%%%%%%%%%%%%%%%%%%
\begin{frame}[fragile]\frametitle{Feynman's Vision: Nature is Quantum}
\begin{itemize}
\item Richard Feynman (1981): ``Nature isn't classical, dammit, and if you want to make a simulation of nature, you'd better make it quantum mechanical''
\item Classical computers struggle to simulate quantum systems: state space is exponential
\item Simulating 50 electrons classically: $2^{50}$ amplitudes = petabytes of memory
\item A quantum computer of the same size can simulate such systems naturally
\item This is arguably the most impactful near-term application of quantum computers
\item Chemistry, materials, pharmaceuticals, clean energy --- all depend on quantum simulation
\end{itemize}
% Suggested diagram: Classical computer struggling vs quantum computer handling molecular simulation
\end{frame}

%%%%%%%%%%%%%%%%%%%%%%%%%%%%%%%%%%%%%%%%%%%%%%%%%%%%%%%%%%%
\begin{frame}[fragile]\frametitle{Why Classical Molecular Simulation is Hard}
\begin{itemize}
\item Molecules obey the Schr\"{o}dinger equation: a quantum mechanical equation
\item Solving exactly for $N$ electrons: exponentially scaling Hilbert space $2^N$ (or larger)
\item H$_2$ (2 electrons): trivial classically; H$_2$O (10 electrons): manageable
\item Penicillin (41 electrons): frontier of classical simulation
\item FeMo-co (nitrogen fixation catalyst, 54 electrons): beyond classical exact solution
\item Approximations (DFT, coupled cluster): fast but imprecise for strongly correlated systems
\end{itemize}
% Suggested diagram: Molecule size vs simulation difficulty (number of electrons vs classical memory)
\end{frame}

%%%%%%%%%%%%%%%%%%%%%%%%%%%%%%%%%%%%%%%%%%%%%%%%%%%%%%%%%%%
\begin{frame}[fragile]\frametitle{VQE for Quantum Chemistry: The Setup}
\begin{itemize}
\item Encode the molecular Hamiltonian as a sum of Pauli operators (qubit operators)
\item Each term in the sum corresponds to electron interactions
\item Use Jordan-Wigner or Bravyi-Kitaev mapping to convert fermionic operators to qubits
\item VQE finds the ground state (lowest energy) of the molecule
\item Ground state energy determines molecular stability, reaction rates, bond lengths
\item Available in Qiskit Nature: built-in molecular Hamiltonians
\end{itemize}
% Suggested diagram: Molecule -> Hamiltonian -> Pauli operators -> quantum circuit pipeline
\end{frame}

%%%%%%%%%%%%%%%%%%%%%%%%%%%%%%%%%%%%%%%%%%%%%%%%%%%%%%%%%%%
\begin{frame}[fragile]\frametitle{Lab: H2 Molecule VQE in Qiskit}
\begin{itemize}
\item Simulate the H$_2$ molecule ground state energy as a function of bond length
\item Use Qiskit Nature to generate the qubit Hamiltonian
\item Run VQE with hardware-efficient ansatz
\item Plot energy vs bond length curve (potential energy surface)
\item Compare with classical exact diagonalization result
\end{itemize}
\begin{verbatim}
from qiskit_nature.second_q.drivers import PySCFDriver
from qiskit_nature.second_q.transformers import FreezeCoreTransformer
driver = PySCFDriver(atom='H 0 0 0; H 0 0 0.735', basis='sto3g')
\end{verbatim}
% Suggested diagram: H2 potential energy surface; VQE vs exact results plot
\end{frame}

%%%%%%%%%%%%%%%%%%%%%%%%%%%%%%%%%%%%%%%%%%%%%%%%%%%%%%%%%%%
\begin{frame}[fragile]\frametitle{Chemical Accuracy and NISQ Limitations}
\begin{itemize}
\item Chemical accuracy: within 1.6 milli-Hartree ($\sim$1 kcal/mol) of exact energy
\item Required for predicting reaction outcomes, binding energies
\item NISQ VQE on small molecules: can approach chemical accuracy with error mitigation
\item Larger molecules: noise prevents accurate results on current hardware
\item Hamiltonian simulation (Trotterization): exact simulation requires deep circuits; not NISQ-friendly
\item Path forward: error mitigation + better ansatze + more qubits
\end{itemize}
\end{frame}

%%%%%%%%%%%%%%%%%%%%%%%%%%%%%%%%%%%%%%%%%%%%%%%%%%%%%%%%%%%
\begin{frame}[fragile]\frametitle{Material Science: Beyond Molecules}
\begin{itemize}
\item Quantum simulation extends to condensed matter systems (crystals, superconductors)
\item Strongly correlated electron systems: classical methods fail, quantum needed
\item High-temperature superconductors: understanding mechanism could revolutionize energy
\item Topological materials: quantum simulators can probe exotic phases of matter
\item Battery materials: optimizing lithium-ion intercalation for better batteries
\item Catalysis: designing industrial catalysts for Haber-Bosch, CO$_2$ reduction
\end{itemize}
% Suggested diagram: Crystal lattice; band structure; quantum phase diagram examples
\end{frame}

%%%%%%%%%%%%%%%%%%%%%%%%%%%%%%%%%%%%%%%%%%%%%%%%%%%%%%%%%%%
\begin{frame}[fragile]\frametitle{Drug Discovery Applications}
\begin{itemize}
\item Drug binding: a drug molecule docks into a protein active site
\item Binding affinity depends on quantum mechanical interactions
\item Classical: docking simulations are fast but use empirical force fields (approximate)
\item Quantum: simulate the exact quantum interaction energy
\item Target: FeMo-co cofactor in nitrogenase (54 electrons) -- needs ~200+ logical qubits
\item IBM estimate: 200-300 qubit fault-tolerant quantum computer could model FeMo-co
\item Real use case: Pfizer, Roche exploring quantum chemistry for drug lead optimization
\end{itemize}
% Suggested diagram: Drug-protein docking visualization; binding energy calculation pipeline
\end{frame}

%%%%%%%%%%%%%%%%%%%%%%%%%%%%%%%%%%%%%%%%%%%%%%%%%%%%%%%%%%%
\begin{frame}[fragile]\frametitle{Financial Quantum Simulation}
\begin{itemize}
\item Monte Carlo simulation: used for option pricing, risk assessment in finance
\item Classical Monte Carlo: $O(1/\epsilon^2)$ samples for error $\epsilon$
\item Quantum amplitude estimation: quadratic speedup for Monte Carlo ($O(1/\epsilon)$ queries)
\item Application: derivative pricing, Value at Risk (VaR) calculations
\item Portfolio optimization: QAOA and VQE formulations developed
\item JP Morgan, Goldman Sachs actively researching quantum finance algorithms
\item Near-term: research; medium-term: potential practical advantage
\end{itemize}
% Suggested diagram: Option price distribution; Monte Carlo sampling vs quantum amplitude estimation
\end{frame}

%%%%%%%%%%%%%%%%%%%%%%%%%%%%%%%%%%%%%%%%%%%%%%%%%%%%%%%%%%%
\begin{frame}[fragile]\frametitle{Hamiltonian Simulation: The General Approach}
\begin{itemize}
\item General goal: simulate time evolution under a Hamiltonian $H$: $e^{-iHt}$
\item Trotterization: split $H$ into simple parts; apply each in short time steps
\item Trotter error: approximation error from non-commuting terms; decreases with step size
\item Qubitization and QSVT: more advanced methods with lower gate counts
\item Hamiltonian simulation is the basis of VQE and quantum chemistry
\item Also used in quantum machine learning (quantum principal component analysis)
\end{itemize}
% Suggested diagram: Trotterization steps; circuit decomposition for exp(iHt)
\end{frame}

%%%%%%%%%%%%%%%%%%%%%%%%%%%%%%%%%%%%%%%%%%%%%%%%%%%%%%%%%%%
\begin{frame}[fragile]\frametitle{Near-Term Quantum Advantage in Simulation}
\begin{itemize}
\item Classical simulation of random quantum circuits: Google's 2019 claim (53 qubits)
\item Counterarguments: classical algorithms improved significantly after the claim
\item Practical quantum advantage in chemistry: likely 100-300 logical qubit range
\item Near-term demonstrations: VQE for small molecules; Trotterized dynamics
\item Quantum Monte Carlo (analog simulation): D-Wave, neutral atom experiments
\item Honest assessment: clear advantage for specific simulation tasks is still being established
\end{itemize}
\end{frame}

%%%%%%%%%%%%%%%%%%%%%%%%%%%%%%%%%%%%%%%%%%%%%%%%%%%%%%%%%%%
\begin{frame}[fragile]\frametitle{Module 9 Summary}
\begin{itemize}
\item Quantum simulation is Feynman's original vision for quantum computing
\item Exponential classical complexity motivates quantum solutions for chemistry/materials
\item VQE: hybrid algorithm for finding molecular ground state energies
\item H$_2$ molecule: demonstrator problem; LiH, BeH$_2$ push current hardware limits
\item Applications: drug discovery, materials design, catalyst optimization, finance
\item NISQ limitations: noise limits accuracy for larger molecules
\item Fault-tolerant QC with 200-300 logical qubits could unlock real chemistry problems
\end{itemize}
\end{frame}
